% ch1_maritime_datasets.tex — Maritime Object Detection Datasets and the Annotation Gap
\section{Maritime Object Detection Datasets and the Annotation Gap}
\label{sec:maritime_datasets}

The performance of modern object detection systems is fundamentally constrained by the
scale, diversity, and quality of their training data. In the autonomous driving domain, this
relationship has been demonstrated repeatedly: each successive generation of benchmark
datasets---from KITTI \cite{Geiger2012KITTI} to nuScenes \cite{Caesar2020nuScenes} to
the Waymo Open Dataset \cite{Sun2020Waymo}---has enabled corresponding advances in
perception algorithms. Maritime autonomy, however, has not benefited from an equivalent
data ecosystem. This chapter surveys the landscape of maritime object detection datasets,
catalogs their scope and annotation methodology, and quantifies the order-of-magnitude gap
that separates maritime data resources from their automotive counterparts. Understanding
this gap is essential context for the broader literature review that follows, which examines
radar-camera calibration, cross-modal projection, automated dataset construction, sensor
fusion, X-band signal processing, and edge deployment---all components of a radar-guided
camera labeling pipeline designed to close the maritime annotation deficit.

\subsection{The Role of Large-Scale Annotated Datasets in Perception}
\label{subsec:role_of_datasets}

The deep learning revolution in computer vision was enabled not only by advances in
network architectures and computing hardware but also by the availability of large-scale,
carefully annotated datasets. ImageNet \cite{Russakovsky2015ImageNet}, with its 1.28
million training images spanning 1,000 object categories, demonstrated that supervised
pre-training on a sufficiently large corpus could yield feature representations that
transferred effectively to downstream tasks. The PASCAL VOC challenge
\cite{everingham2010pascal} established standardized evaluation protocols that allowed
meaningful comparison across detection methods, while MS~COCO \cite{Lin2014COCO}
raised the bar with 330,000 images, 1.5 million object instances, and dense annotations
including bounding boxes, segmentation masks, and captions. In the autonomous driving
domain, this progression continued: KITTI \cite{Geiger2012KITTI} introduced multi-sensor
benchmarks with LiDAR, stereo cameras, and GPS/IMU; nuScenes \cite{Caesar2020nuScenes}
scaled to 1,000 scenes with full 360-degree sensor coverage and 1.4 million 3D bounding
box annotations; and the Waymo Open Dataset \cite{Sun2020Waymo} reached 12.6 million
LiDAR box labels and 11.8 million camera box labels across 1,150 driving scenes.

The impact of dataset scale on detector performance is well documented. Zhang et
al.~\cite{Zhang2021MarineDetectionSurvey} note that large-scale, multi-scenario training
data is an indispensable prerequisite for practical deployment of deep learning-based
detection systems, and that the maritime domain lacks a widely accepted standardized
benchmark of comparable scale. This observation motivates a systematic examination of
what maritime datasets do exist, what they contain, and where the critical deficiencies lie.

\subsection{A Catalog of Maritime Object Detection Datasets}
\label{subsec:dataset_catalog}

Maritime detection datasets have emerged from three distinct operational contexts:
shore-based surveillance systems monitoring harbors and coastal waters, onboard sensors
carried by unmanned surface vehicles (USVs) or autonomous surface vehicles (ASVs), and
multi-modal platforms integrating radar, LiDAR, or other active sensors alongside cameras.
A fourth category encompasses specialized datasets that employ infrared, hyperspectral,
or stereo imaging. The following subsections organize the literature accordingly.

\subsubsection{Shore-Based Visual Datasets}
\label{subsubsec:shore_based}

Shore-based maritime datasets are typically collected from fixed camera installations
overlooking harbors, canals, or coastal waterways. Their primary applications are
maritime surveillance, vessel traffic monitoring, and port security.

\textbf{MarDCT} (2015). The Maritime Detection, Classification, and Tracking benchmark
\cite{Bloisi2015MarDCT} was among the earliest publicly available maritime video datasets.
Collected by the ARGOS automated surveillance system operating continuously in Venice,
Italy, MarDCT comprises 6,742 images and 16 video sequences at $800 \times 240$
resolution. Ground truth is provided as bounding boxes for boat detection. While
pioneering in its public availability, the dataset's small scale and limited resolution
restrict its utility for training modern deep detectors.

\textbf{Singapore Maritime Dataset (SMD)} (2017). Prasad et al.~\cite{Prasad2017SMD}
released the Singapore Maritime Dataset alongside a comprehensive survey of video
processing for maritime object detection. The dataset contains 81 high-definition
($1920 \times 1080$) video sequences captured from both onshore (32 sequences) and
onboard (4 sequences) viewpoints in Singapore waters, using visual-optical and
near-infrared sensors. Annotations cover 10 maritime object classes across diverse
environmental conditions including haze, rain, sunrise, midday, and nighttime. SMD
established an important early benchmark for evaluating detection and tracking algorithms
in tropical maritime environments.

\textbf{SeaShips} (2018). Shao et al.~\cite{Shao2018SeaShips} constructed a large-scale,
precisely annotated dataset from approximately 10,080 real-world video segments captured
by a deployed coastline surveillance system. SeaShips contains 31,455 images covering six
common vessel types: ore carrier, bulk cargo carrier, general cargo ship, container ship,
fishing boat, and passenger ship. Images were carefully selected to represent variations
in scale, hull visibility, illumination, viewpoint, background complexity, and occlusion.
Each image is annotated with ship type labels and high-precision bounding boxes. At the
time of its release, SeaShips was among the largest publicly available maritime detection
datasets and remains widely used as a benchmark.

\textbf{MCShips} (2020). Zheng and Zhang \cite{Zheng2020MCShips} introduced MCShips, a
dataset of 14,709 annotated images collected via web crawlers and search engines,
spanning 13 fine-grained ship categories (6 warship types and 7 civilian vessel types).
Unlike surveillance-derived datasets, MCShips captures a wide diversity of imaging
conditions, perspectives, and resolutions found in publicly available online imagery.
Each image is annotated with bounding boxes and fine-grained class labels, supporting
both detection and classification tasks.

\textbf{ABOShips} (2021). Iancu et al.~\cite{Iancu2021ABOShips} constructed the ABOShips
dataset from video sequences recorded in the Finnish Archipelago using a camera mounted
on a waterbus. The dataset comprises 9,880 images containing 41,967 annotated object
instances across 9 vessel categories plus seamarks and miscellaneous floaters. A two-stage
annotation process was employed: initial manual labeling followed by automated tracking
(CSRT tracker) to identify and correct labeling inconsistencies. Benchmark evaluations
with Faster R-CNN, R-FCN, SSD, and EfficientDet established baseline performance levels,
with Faster R-CNN (Inception-ResNet v2 backbone) achieving the highest overall accuracy.

\textbf{MCMOD} (2023). Sun et al.~\cite{Sun2023MCMOD} released the Multi-Category
Maritime Object Detection dataset, consisting of 16,166 images with 98,590 annotated
objects across 10 maritime categories including fishing vessels, speedboats, engineering
ships, cargo ships, yachts, sailboats, buoys, rafts, cruise ships, and an ``others''
class. Three onshore cameras operating continuously in Hainan, China, captured
high-resolution ($1920 \times 1080$) video under varied conditions including fog, rain,
and evening illumination. The cameras' rotation and zoom capabilities produced diverse
scenes, and the dataset's emphasis on small-scale distant targets addresses a persistent
challenge in maritime detection.

\textbf{KOLOMVERSE} (2024). Nanda et al.~\cite{Nanda2024KOLOMVERSE} introduced the
Korea Open Large-Scale Image Dataset for Object Detection in the Maritime Universe,
representing the largest publicly available maritime detection dataset to date. Drawn
from 5,845 hours of video data captured across 21 territorial waters of South Korea,
KOLOMVERSE contains 186,419 4K-resolution ($3840 \times 2160$) images with 732,039
annotated object instances across five classes: ship (393,936), buoy (34,080), fishnet
buoy (95,815), lighthouse (60,362), and wind farm (147,846). The dataset is partitioned
into training (149,175), validation (18,643), and test (18,601) splits. KOLOMVERSE's
scale and resolution represent a significant step forward for maritime detection,
though its five-class taxonomy and shore-based perspective limit its applicability to
fine-grained classification and onboard perception tasks.

\subsubsection{USV/ASV-Mounted Datasets}
\label{subsubsec:usv_mounted}

Datasets collected from the perspective of an unmanned or autonomous surface vehicle
present unique challenges not captured by shore-based imagery: low camera elevation
above the waterline, dynamic horizon orientation due to wave-induced platform motion,
frequent spray and lens contamination, and the need to detect obstacles at close range
for collision avoidance.

\textbf{MaSTr1325} (2019). Bovcon et al.~\cite{Bovcon2019MaSTr1325} released a maritime
semantic segmentation training dataset containing 1,325 diverse images captured over two
years by a USV operating in the Gulf of Koper, Slovenia. Each image is annotated at
pixel level with three semantic classes: sea, sky, and environment (encompassing all
obstacles and shoreline). While the three-class taxonomy is coarse, the dataset was
instrumental in training early deep segmentation models for USV obstacle detection. The
annotation procedure required approximately twenty minutes per image, underscoring the
labor intensity of pixel-level maritime labeling.

\textbf{FloW} (2021). Cheng et al.~\cite{Cheng2021FloW} presented FloW, the first
dataset targeting floating waste detection in inland waterways, accepted at ICCV 2021.
FloW comprises two sub-datasets: FloW-Img, an image collection with bounding box
annotations for floating debris, and FloW-RI, a multimodal sub-dataset containing
synchronized millimeter-wave radar data and camera images. The inclusion of radar data
makes FloW one of the earliest maritime-adjacent datasets to provide cross-modal
supervision signals. Baseline experiments demonstrated that floating waste detection
remains challenging due to small target size, strong water-surface reflections, and
bankside clutter.

\textbf{MassMIND} (2023). Nirgudkar et al.~\cite{Nirgudkar2023MassMIND} contributed
the Massachusetts Maritime Infrared Dataset, containing 2,916 long-wave infrared (LWIR)
images captured over two years in and around Boston Harbor. The dataset is segmented
using both instance and semantic segmentation into seven classes, making it the first
maritime dataset to provide fine-grained LWIR segmentation labels for both living
(humans, marine wildlife) and non-living (vessels, structures) obstacles. Evaluation with
UNet, PSPNet, and DeepLabv3 architectures established baseline segmentation performance
for thermal maritime imagery.

\textbf{LaRS} (2023). \v{Z}ust et al.~\cite{Zust2023LaRS} released the first maritime
panoptic obstacle detection benchmark, accepted at ICCV 2023. LaRS contains over 4,000
per-pixel labeled keyframes with nine preceding temporal context frames each, totaling
more than 40,000 frames. Annotations span 8 ``thing'' classes, 3 ``stuff'' classes, and
19 global scene attributes. The dataset features the largest diversity among maritime
datasets in recording locations (lakes, rivers, and seas), scene types, obstacle classes,
and acquisition conditions. An online evaluation server with results from 27 segmentation
methods provides a standardized benchmark. LaRS represents the current state of the art
for USV-perspective semantic and panoptic segmentation evaluation.

\textbf{WaterScenes} (2024). Zhang et al.~\cite{Zhang2024WaterScenes} introduced
WaterScenes, the first multi-task 4D radar-camera fusion dataset for autonomous driving
on water surfaces. Collected by a USV equipped with a 4D millimeter-wave radar and a
monocular camera, the dataset comprises 54,120 synchronized image-radar frame pairs
covering more than 200,000 annotated objects. Annotations support five tasks: object
detection, instance segmentation, semantic segmentation, free-space segmentation, and
waterline segmentation. Environmental diversity spans daytime, nightfall, and nighttime
conditions; normal, dim, and strong lighting; sunny, overcast, rainy, and snowy weather;
and four waterway types (river, lake, canal, moat). WaterScenes is notable for being the
first maritime dataset to provide both pixel-level camera annotations and point-level
radar annotations, enabling direct evaluation of radar-camera fusion algorithms in the
maritime domain.

\textbf{SeePerSea} (2024). Jeong et al.~\cite{Jeong2024SeePerSea} released the first
publicly accessible multi-modal perception dataset for in-water obstacle detection from
ASVs, featuring synchronized LiDAR point clouds and RGB images. Collected over four years
across locations in the United States, Barbados, and South Korea, SeePerSea contains
11,561 annotated frames with navigation-oriented three-class labels (ship, buoy, other).
The dataset is significant as the first maritime LiDAR-camera dataset with object labels
across both modalities, supporting evaluation of cross-modal fusion for aquatic obstacle
detection.

\subsubsection{Multi-Modal Radar/LiDAR/Camera Datasets}
\label{subsubsec:multimodal}

The most recent generation of maritime datasets integrates multiple active and passive
sensing modalities, reflecting the recognition that robust maritime perception requires
sensor fusion. These datasets are also the most directly relevant to radar-guided camera
labeling, as they provide synchronized radar and camera data that can serve as ground
truth for evaluating cross-modal projection and automatic annotation pipelines.

\textbf{Pohang Canal Dataset} (2023). Chung et al.~\cite{Chung2023Pohang} released a
comprehensive multi-modal maritime dataset collected along a 7.5~km route through the
Pohang Canal, inner and outer ports, and near-coastal areas in South Korea. The sensor
suite includes stereo cameras, an infrared camera, an omnidirectional camera, three
LiDARs, a marine radar, GPS, and an attitude heading reference system. Data were collected
under diverse weather and visibility conditions. The Pohang Canal Dataset is among the
most sensor-rich maritime collections available, though it is primarily designed for
navigation and SLAM rather than object detection, and its annotation density for detection
tasks is limited.

\textbf{M$^2$SODAI} (2023). Jang et al.~\cite{Jang2023M2SODAI} presented a multi-modal
maritime object detection dataset combining RGB and hyperspectral image sensors, accepted
at NeurIPS 2023 Datasets and Benchmarks track. The dataset provides synchronized aerial
image pairs with bounding box annotations across multiple maritime object categories.
Hyperspectral imaging (HSI), which captures over 100 spectral channels in the visible
and near-infrared range, enables extraction of material-specific signatures that can
discriminate maritime objects from sea-surface clutter even when visual appearance is
ambiguous. M$^2$SODAI is the first maritime dataset to incorporate hyperspectral sensing,
addressing the false-positive challenge inherent in single-modality visual detection
over dynamic water surfaces.

\textbf{PoLaRIS} (2024). Choi et al.~\cite{Choi2024PoLaRIS} released the PoLaRIS
(Pohang Canal Object Detection and Tracking Dataset in Maritime Environments) dataset,
accepted at ICRA 2025. PoLaRIS provides approximately 360,000 labeled images from RGB
and thermal infrared cameras, along with point-wise annotations from radar and LiDAR. It
is the first dataset offering multi-modal annotations specifically tailored to maritime
object detection and tracking, including tracking ground truth for dynamic objects with
bounding boxes projected across modalities. The dataset includes annotations for objects
as small as $10 \times 10$ pixels, reflecting the prevalence of small, distant targets
in maritime environments. With approximately 190,000 detection-grade bounding box
annotations, PoLaRIS represents the most densely annotated multi-sensor maritime
dataset available.

\textbf{MOANA} (2024). Jang et al.~\cite{Jang2024MOANA} introduced the Multi-Radar
Dataset for Maritime Odometry and Autonomous Navigation Application, integrating
short-range LiDAR, medium-range W-band radar, and long-range X-band radar into a unified
multi-scale sensing framework. The dataset comprises 7 sequences collected from diverse
coastal regions with varying navigation difficulty. MOANA is the first maritime dataset
to combine multiple radar frequency bands, enabling research on multi-scale maritime
perception from close-range obstacle detection to long-range situational awareness. Its
primary focus is odometry estimation, SLAM, and place recognition rather than object
detection, and the limited number of sequences constrains its statistical power for
training data-hungry detection models.

\subsubsection{Multi-Spectral and Specialized Datasets}
\label{subsubsec:specialized}

Several datasets address maritime perception through non-standard imaging modalities or
specialized acquisition configurations.

\textbf{VAIS} (2015). Zhang et al.~\cite{Zhang2015VAIS} released the Visible and
Infrared Ship dataset, the first publicly available collection of paired visible and
thermal infrared maritime imagery. VAIS contains over 1,000 paired RGB-IR images across
six vessel categories (merchant, sailing, passenger, medium-other, tug, and small), all
acquired from pier-side vantage points. The simultaneous dual-band acquisition enables
research on cross-spectral recognition and day-night vessel classification, though the
dataset's modest scale limits its applicability for training deep networks.

\textbf{MARVEL} (2016). Gundogdu et al.~\cite{Gundogdu2016MARVEL} constructed a
large-scale maritime vessel image dataset by harvesting approximately 2 million
user-uploaded images and metadata from an online ship-spotting community. Images are
categorized into 26 superclasses constructed from 109 fine-grained vessel type labels,
with additional attributes including vessel identity, photograph category, and year of
construction. MARVEL supports vessel classification, verification, retrieval, and
recognition tasks. However, the dataset's web-sourced provenance means images are
captured under uncontrolled conditions without standardized annotations, and no bounding
box or segmentation labels are provided, limiting its utility for detection tasks.

\textbf{MODS} (2022). Bovcon et al.~\cite{Bovcon2022MODS} released a USV-oriented object
detection and obstacle segmentation benchmark containing approximately 81,000 stereo
images synchronized with onboard IMU measurements, with over 60,000 annotated objects.
MODS addresses two major perception tasks---maritime object detection and maritime
obstacle segmentation---and introduces a standardized evaluation protocol that measures
detection accuracy in terms meaningful for practical USV navigation. Benchmark results
from 19 state-of-the-art detection and segmentation methods provide a comprehensive
performance landscape for USV perception.

\subsection{Comprehensive Dataset Comparison}
\label{subsec:dataset_comparison}

Table~\ref{tab:maritime_datasets} presents a systematic comparison of the maritime
datasets surveyed in this chapter. The table reveals several patterns: (1) dataset scale
has grown substantially over the past decade, from thousands to hundreds of thousands of
images; (2) the majority of datasets remain single-modality visual collections; (3)
multi-modal datasets with radar or LiDAR are a recent development (2023--2024) and remain
small by comparison; and (4) annotation granularity varies widely, from image-level class
labels to pixel-level panoptic segmentation.

\begin{table}[!htbp]
\centering
\caption{Comprehensive comparison of maritime object detection datasets. Annotation types:
BBox = bounding box, Seg = semantic segmentation, Inst = instance segmentation,
Pan = panoptic segmentation, PW = point-wise, Trk = tracking ID.}
\label{tab:maritime_datasets}
\small
\begin{tabularx}{\textwidth}{@{}l c l r r c X@{}}
\toprule
\textbf{Dataset} & \textbf{Year} & \textbf{Sensors} & \textbf{Images} & \textbf{Objects} & \textbf{Classes} & \textbf{Annotation} \\
\midrule
MarDCT \cite{Bloisi2015MarDCT}         & 2015 & RGB            &      6,742  &      ---   & Boat     & BBox \\
VAIS \cite{Zhang2015VAIS}              & 2015 & RGB + IR       &      2,000  &    1,242   & 6        & BBox \\
MARVEL \cite{Gundogdu2016MARVEL}       & 2016 & RGB (web)      &  2,000,000  &      ---   & 26       & Class label \\
SMD \cite{Prasad2017SMD}               & 2017 & RGB + NIR      &     81 vid  &      ---   & 10       & BBox (video) \\
SeaShips \cite{Shao2018SeaShips}       & 2018 & RGB            &     31,455  &   40,077   & 6        & BBox \\
MaSTr1325 \cite{Bovcon2019MaSTr1325}   & 2019 & RGB + IMU      &      1,325  &      ---   & 3        & Seg (pixel) \\
MCShips \cite{Zheng2020MCShips}        & 2020 & RGB (web)      &     14,709  &   14,709   & 13       & BBox \\
ABOShips \cite{Iancu2021ABOShips}      & 2021 & RGB            &      9,880  &   41,967   & 9+       & BBox \\
FloW \cite{Cheng2021FloW}              & 2021 & RGB + mmW radar &    2,000+  &    6,000+  & Waste    & BBox + PW \\
MODS \cite{Bovcon2022MODS}             & 2022 & Stereo + IMU   &     81,000  &   60,000+  & Multi    & BBox + Seg \\
MassMIND \cite{Nirgudkar2023MassMIND}  & 2023 & LWIR           &      2,916  &      ---   & 7        & Seg + Inst \\
LaRS \cite{Zust2023LaRS}               & 2023 & RGB            &     40,000+ &      ---   & 8+3      & Pan (pixel) \\
M$^2$SODAI \cite{Jang2023M2SODAI}      & 2023 & RGB + HSI      &      5,764+ &      ---   & Multi    & BBox \\
Pohang \cite{Chung2023Pohang}          & 2023 & Multi (7+)     &       ---   &      ---   & ---      & Nav. labels \\
MCMOD \cite{Sun2023MCMOD}              & 2023 & RGB            &     16,166  &   98,590   & 10       & BBox \\
KOLOMVERSE \cite{Nanda2024KOLOMVERSE}  & 2024 & RGB            &    186,419  &  732,039   & 5        & BBox \\
WaterScenes \cite{Zhang2024WaterScenes}& 2024 & RGB + 4D radar &     54,120  &  200,000+  & Multi    & BBox + Seg + PW \\
PoLaRIS \cite{Choi2024PoLaRIS}        & 2024 & RGB+IR+LiDAR+R &    360,000  & $\sim$190k & Multi    & BBox + PW + Trk \\
SeePerSea \cite{Jeong2024SeePerSea}    & 2024 & RGB + LiDAR    &     11,561  &      ---   & 3        & BBox + PW \\
MOANA \cite{Jang2024MOANA}             & 2024 & Multi-radar+L  &    7 seq.   &      ---   & ---      & Nav. labels \\
\bottomrule
\end{tabularx}
\end{table}

Several observations emerge from this comparison. First, KOLOMVERSE \cite{Nanda2024KOLOMVERSE}
is the only maritime dataset to exceed 100,000 detection-annotated images, yet it covers
only five object classes from a shore-based perspective. Second, WaterScenes
\cite{Zhang2024WaterScenes} is the only dataset providing both dense camera annotations
and radar point annotations from a vessel-mounted platform, making it uniquely relevant
for radar-camera fusion research. Third, the multi-sensor datasets most suitable for
evaluating radar-guided labeling---Pohang \cite{Chung2023Pohang}, PoLaRIS
\cite{Choi2024PoLaRIS}, and MOANA \cite{Jang2024MOANA}---either lack dense detection
annotations or comprise only a handful of sequences.

\subsection{The Automotive Benchmark: Scale and Infrastructure}
\label{subsec:automotive_comparison}

To contextualize the maritime annotation gap, it is instructive to examine the scale of
automotive perception benchmarks that have driven progress in self-driving vehicle
perception. Table~\ref{tab:automotive_datasets} summarizes the key automotive datasets.

\begin{table}[!htbp]
\centering
\caption{Major automotive perception benchmark datasets for comparison with maritime
equivalents.}
\label{tab:automotive_datasets}
\small
\begin{tabularx}{\textwidth}{@{}l c l r r c X@{}}
\toprule
\textbf{Dataset} & \textbf{Year} & \textbf{Sensors} & \textbf{Images} & \textbf{Labels} & \textbf{Classes} & \textbf{Annotation} \\
\midrule
ImageNet \cite{Russakovsky2015ImageNet}  & 2012 & RGB          & 1,281,167 &  ---           & 1,000 & Class label \\
KITTI \cite{Geiger2012KITTI}             & 2012 & Stereo+LiDAR &    14,999 &  200,000+      & 8     & 3D BBox \\
MS~COCO \cite{Lin2014COCO}              & 2014 & RGB          &   330,000 &  1,500,000     & 80    & BBox + Seg \\
nuScenes \cite{Caesar2020nuScenes}       & 2020 & 6C+5R+1L     & 1,400,000 &  1,400,000     & 23    & 3D BBox + Trk \\
Waymo \cite{Sun2020Waymo}                & 2020 & 5C+5L        & ---       &  12,600,000    & 4     & 3D BBox + Trk \\
\bottomrule
\end{tabularx}
\end{table}

The automotive ecosystem benefits from a virtuous cycle: large datasets enable better
models, which attract more investment, which funds further data collection and annotation.
The Waymo Open Dataset alone contains 12.6 million LiDAR 3D bounding box labels with
tracking identifiers across 1,150 scenes, each spanning 20 seconds of driving with
well-synchronized and calibrated LiDAR and camera data \cite{Sun2020Waymo}. nuScenes
provides 1.4 million 3D bounding box annotations across 1,000 scenes with full
360-degree sensor coverage including six cameras, five radars, and one LiDAR
\cite{Caesar2020nuScenes}. The automotive radar-camera datasets RADIATE
\cite{sheeny2020radiate} and the Oxford Radar RobotCar Dataset \cite{barnes2020oxford}
further demonstrate the automotive community's investment in multi-modal data collection.

\subsection{Quantifying the Maritime--Automotive Gap}
\label{subsec:gap_analysis}

Table~\ref{tab:gap_analysis} presents a direct quantitative comparison between the
best-resourced maritime and automotive datasets across several key metrics. The gaps are
stark and consistent across all dimensions.

\begin{table}[!htbp]
\centering
\caption{Order-of-magnitude gap between automotive and maritime dataset resources across
key metrics. Gap factor is computed as the ratio of automotive to maritime values.}
\label{tab:gap_analysis}
\small
\begin{tabularx}{\textwidth}{@{}X r r r@{}}
\toprule
\textbf{Metric} & \textbf{Automotive Best} & \textbf{Maritime Best} & \textbf{Gap Factor} \\
\midrule
Detection images (single dataset)    & 1,281,167 (ImageNet)  & 186,419 (KOLOMVERSE) & $\sim$7$\times$ \\
3D bbox labels (single dataset)      & 12,600,000 (Waymo)    & $\sim$190,000 (PoLaRIS) & $\sim$66$\times$ \\
Multi-sensor scenes                  & 2,150 (Waymo+nuScenes) & 7 (MOANA)           & $\sim$307$\times$ \\
Annotation instances (total)         & $>$15,000,000         & $\sim$1,300,000       & $\sim$12$\times$ \\
Sensor modalities per dataset (max)  & 12 (nuScenes)         & 7+ (Pohang/PoLaRIS)  & $\sim$2$\times$ \\
Detection classes (max, single)      & 80 (COCO)             & 13 (MCShips)         & $\sim$6$\times$ \\
\bottomrule
\end{tabularx}
\end{table}

The most critical gap is in multi-sensor data volume. The Waymo Open Dataset provides
1,150 fully annotated multi-sensor scenes, and nuScenes adds another 1,000, for a
combined total exceeding 2,000 scenes with synchronized LiDAR, camera, and radar data.
By comparison, MOANA \cite{Jang2024MOANA}---the most comprehensive multi-radar maritime
dataset---comprises only 7 sequences. Even aggregating all maritime multi-modal datasets
(Pohang, PoLaRIS, MOANA, WaterScenes, FloW, SeePerSea), the total number of
independently collected multi-sensor sequences remains well below 100. This
approximately 300-fold gap in multi-sensor scene diversity represents perhaps the most
consequential deficiency for developing robust maritime perception systems, as sensor
fusion algorithms require diverse training examples to generalize across the wide
range of environmental conditions encountered at sea.

The 3D bounding box label gap is similarly severe. Waymo's 12.6 million LiDAR-derived
3D labels dwarf the approximately 190,000 detection-grade annotations in PoLaRIS
\cite{Choi2024PoLaRIS}, yielding a gap factor of roughly 66. This disparity is
compounded by the fact that maritime 3D labels are almost exclusively derived from
2D bounding boxes (sometimes augmented with range information), whereas automotive
3D labels benefit from dense LiDAR point clouds that enable precise 3D cuboid
fitting.

Aggregating across all maritime datasets surveyed in this chapter, the total number of
detection-grade annotated maritime images (those with bounding box or segmentation
annotations) is approximately 700,000. This figure is dominated by KOLOMVERSE at
186,419, followed by PoLaRIS at approximately 360,000 (though many of these are
repetitive keyframes from overlapping sequences), MODS at 81,000, WaterScenes at
54,120, SeaShips at 31,455, and the remaining datasets collectively contributing
under 100,000. For comparison, MS~COCO alone provides 330,000 annotated images with
1.5 million object instances \cite{Lin2014COCO}, and the Waymo Open Dataset
encompasses over 200,000 annotated LiDAR frames \cite{Sun2020Waymo}.

\subsection{Annotation Cost and Quality Challenges}
\label{subsec:annotation_challenges}

The maritime annotation gap reflects not merely a shortage of collection infrastructure
but also the disproportionate difficulty and cost of annotating maritime imagery compared
to terrestrial scenes.

\subsubsection{Domain Expertise Requirements}

Maritime object annotation demands specialized knowledge that general-purpose crowd
workers typically lack. Distinguishing between vessel types---cargo ships, tankers,
fishing trawlers, naval vessels, recreational craft---requires familiarity with hull
forms, superstructure configurations, and operational context that is not readily
acquired without maritime training. Zhang et al.~\cite{Zhang2021MarineDetectionSurvey}
emphasize that domain-specific expertise is a prerequisite for constructing
high-quality maritime training data, noting that misclassification at the labeling
stage propagates directly into detector training as systematic label noise. The ABOShips
annotation pipeline \cite{Iancu2021ABOShips}, which required a second pass using visual
tracking to identify and correct initial labeling errors, illustrates the additional
quality-control overhead inherent in maritime annotation.

\subsubsection{Small and Distant Targets}

Maritime scenes routinely contain targets at ranges of several kilometers, appearing as
clusters of only a few pixels in the image. MCMOD \cite{Sun2023MCMOD} explicitly
addresses the prevalence of small-scale objects, and PoLaRIS \cite{Choi2024PoLaRIS}
includes annotations for objects as small as $10 \times 10$ pixels. Annotating such
targets is time-consuming and error-prone: annotators must zoom to high magnification
to delineate object boundaries, and the low signal-to-noise ratio for distant targets
increases inter-annotator disagreement. The annotation time per image for maritime
datasets is correspondingly high; Bovcon et al.~\cite{Bovcon2019MaSTr1325} report
approximately 20 minutes per image for pixel-level semantic segmentation of MaSTr1325,
compared to typical rates of 5--10 minutes for road-scene segmentation in automotive
datasets.

\subsubsection{Ambiguous Boundaries and Visual Clutter}

Water-surface conditions create annotation ambiguities that are largely absent from
terrestrial scenes. Vessel wakes, bow waves, and spray obscure hull boundaries; sun
glare and specular reflections from the water surface produce bright artifacts that
can be confused with objects; and atmospheric effects (haze, fog, rain) reduce contrast
between objects and background. These factors increase the variance of bounding box
placement across annotators. Additionally, partially submerged objects, overlapping
vessels at marinas, and the lack of fixed spatial reference points (unlike lane
markings or curbs in driving scenes) further complicate boundary delineation.

\subsubsection{Class Imbalance and Long-Tail Distribution}

Maritime traffic exhibits severe class imbalance. In KOLOMVERSE
\cite{Nanda2024KOLOMVERSE}, ships account for 393,936 of 732,039 instances (53.8\%),
while buoys contribute only 34,080 (4.7\%). Similar distributions appear in other
datasets: small recreational craft dominate in coastal waters, while large commercial
vessels are relatively rare. This long-tail distribution means that even large datasets
may contain too few examples of safety-critical minority classes (e.g., kayaks,
swimmers, semi-submerged debris) to train reliable detectors.

\subsection{Domain-Specific Perception Challenges}
\label{subsec:domain_challenges}

Beyond annotation difficulty, the maritime domain presents perception challenges that
necessitate larger and more diverse training sets than might be required for a
comparably complex terrestrial task.

\subsubsection{Environmental Variability}

Maritime environments exhibit extreme variability across temporal, meteorological, and
geographic dimensions. Sea state ranges from mirror-calm to storm conditions with
multi-meter waves; lighting spans from direct tropical sun (with intense specular
reflections) to overcast polar twilight; visibility ranges from unlimited to near-zero
in fog, rain, or snow. WaterScenes \cite{Zhang2024WaterScenes} explicitly captures
variation across daytime, nightfall, and nighttime conditions; sunny, overcast, rainy,
and snowy weather; and four distinct waterway types. KOLOMVERSE
\cite{Nanda2024KOLOMVERSE} samples from 21 territorial waters to achieve geographic
diversity. Despite these efforts, no single maritime dataset approaches the
environmental coverage that automotive datasets achieve through crowd-sourced
collection across entire nations and seasons.

\subsubsection{Horizon and Background Complexity}

The maritime horizon, while seemingly simple, creates distinctive detection challenges.
Objects transitioning from below to above the horizon change dramatically in
background context (dark water to bright sky or vice versa). Coastal backgrounds
introduce shoreline structures, breakwaters, and moored vessels that create complex
clutter patterns. In narrow waterways, such as those captured in the Pohang Canal
Dataset \cite{Chung2023Pohang}, the visual scene more closely resembles an urban
canyon than open ocean, requiring detectors to handle both maritime and semi-terrestrial
visual contexts.

\subsubsection{Scale and Aspect Ratio Diversity}

The range of object scales in maritime scenes exceeds that in typical driving scenarios.
A single image may contain both a large vessel filling much of the frame and distant
targets comprising fewer than 100 pixels. Moreover, the aspect ratios of maritime
objects vary more widely than those of vehicles, pedestrians, and cyclists: a container
ship viewed broadside may have an aspect ratio exceeding 8:1, while a small buoy is
approximately 1:1. This scale diversity challenges both anchor-based and anchor-free
detection architectures, which must maintain sensitivity across several orders of
magnitude in object area.

\subsubsection{Temporal Dynamics and Motion Patterns}

Wave-induced platform motion on USV-mounted cameras produces rapid and largely
unpredictable changes in camera orientation, leading to motion blur, horizon
oscillation, and intermittent spray occlusion. The MODS dataset \cite{Bovcon2022MODS}
and LaRS \cite{Zust2023LaRS} both capture these dynamics through temporally consecutive
frame sequences. Additionally, the relative motion between an observing platform and
maritime targets differs qualitatively from automotive scenarios: closing rates can be
very slow (vessels approaching head-on at a few knots) or very fast (high-speed craft),
and lateral apparent motion is dominated by the observer's heading changes rather than
lane-based geometry.

\subsection{Summary and Research Gaps}
\label{subsec:ch1_summary}

This chapter has surveyed 20 maritime object detection datasets spanning a decade of
development (2015--2024), organized by collection platform and sensor modality. Several
structural factors account for the persistent annotation gap between maritime and
automotive domains:

\begin{enumerate}[leftmargin=*]
\item \textbf{Limited collection infrastructure.} Automotive datasets benefit from
    fleets of sensor-equipped vehicles operating on public roads; maritime data
    collection requires access to vessels, coastal installations, or USVs that are
    less widely available and more expensive to operate.

\item \textbf{Fewer research groups.} The autonomous driving community encompasses
    dozens of major industrial laboratories (Waymo, Cruise, Argo AI, Mobileye, Baidu
    Apollo) and hundreds of academic groups, each contributing data and benchmarks. The
    maritime perception community is substantially smaller, with fewer organizations
    possessing both the sensing infrastructure and the annotation resources to produce
    large-scale datasets.

\item \textbf{Environmental diversity.} Maritime environments span open ocean, coastal
    waters, rivers, lakes, canals, and harbors, each with distinct visual
    characteristics, traffic patterns, and object taxonomies. Achieving representative
    coverage requires data collection across many geographically and climatically
    distinct locations, a logistical challenge that compounds the per-image annotation
    cost.

\item \textbf{Annotation difficulty.} As documented in
    Section~\ref{subsec:annotation_challenges}, the combination of small distant
    targets, ambiguous boundaries, domain expertise requirements, and severe class
    imbalance makes maritime annotation two to four times slower per image than
    automotive annotation. At typical annotation service rates, the cost of manually
    labeling a maritime dataset to Waymo-equivalent scale (12 million bounding boxes)
    would be prohibitive for most research groups.
\end{enumerate}

Against this backdrop, several emerging approaches offer pathways to narrow the gap:

\begin{enumerate}[leftmargin=*]
\item \textbf{Radar-guided camera labeling.} By projecting radar detections into the
    camera frame, bounding box annotations can be generated automatically without human
    intervention. The availability of WaterScenes \cite{Zhang2024WaterScenes}, PoLaRIS
    \cite{Choi2024PoLaRIS}, and MOANA \cite{Jang2024MOANA} as multi-modal reference
    datasets provides initial validation targets for such pipelines. The noise
    characteristics of radar-projected labels are systematic (arising from calibration
    and projection geometry) rather than random, which suggests that domain-specific
    denoising strategies may be required.

\item \textbf{AIS-assisted annotation.} The Automatic Identification System (AIS)
    provides vessel identity, position, heading, and speed for cooperative maritime
    traffic. Fusing AIS tracks with camera imagery enables automatic association of
    detected objects with vessel metadata, facilitating both bounding box generation
    (via projected AIS positions) and fine-grained class labeling (via vessel type
    codes transmitted in AIS messages).

\item \textbf{Active learning and selective annotation.} Rather than annotating all
    collected frames uniformly, active learning strategies can identify the most
    informative frames for human annotation, maximizing detection improvement per
    annotation dollar. The temporal redundancy inherent in maritime video (vessels
    move slowly relative to frame rate) suggests that aggressive frame selection
    could substantially reduce annotation volume without sacrificing model performance.

\item \textbf{Cross-modal label transfer.} Beyond radar-to-camera projection,
    cross-modal transfer encompasses LiDAR-to-camera (as explored in SeePerSea
    \cite{Jeong2024SeePerSea}), thermal-to-visible (as enabled by VAIS
    \cite{Zhang2015VAIS} and PoLaRIS \cite{Choi2024PoLaRIS}), and
    hyperspectral-to-RGB (as proposed in M$^2$SODAI \cite{Jang2023M2SODAI}).
    Each modality pair offers distinct advantages: radar provides range and velocity;
    LiDAR provides precise 3D geometry; thermal provides illumination-invariant
    detection; and hyperspectral provides material discrimination. A comprehensive
    automated labeling framework may ultimately combine multiple cross-modal signals
    to produce annotations that exceed the quality achievable by any single modality
    alone.
\end{enumerate}

The remainder of this literature review examines each of these technical components in
detail, beginning with radar-camera sensor registration and calibration
(Section~\ref{sec:radar_camera_calibration}), progressing through radar-to-camera
projection and bounding box generation (Section~\ref{sec:projection_bboxgen}), and
extending to automated dataset construction methodologies
(Section~\ref{sec:automated_datasets}), radar-camera fusion
(Section~\ref{sec:radar_camera_fusion}), X-band signal processing
(Section~\ref{sec:xband_processing}), and edge computing for real-time deployment
(Section~\ref{sec:edge_computing}).
